\documentclass[11pt]{article}
\usepackage[margin=1in]{geometry}
\usepackage{amsmath,amssymb,amsthm}
\usepackage[colorlinks=true,urlcolor=blue,linkcolor=blue,citecolor=blue]{hyperref}
\newtheorem{theorem}{Theorem}[section]
\newtheorem{lemma}[theorem]{Lemma}
\newtheorem{corollary}[theorem]{Corollary}
\newtheorem{remark}[theorem]{Remark}
\newcommand{\diag}{\operatorname{diag}}
\newcommand{\spec}{\operatorname{spec}}
\newcommand{\rank}{\operatorname{rank}}
\newcommand{\norm}[1]{\left\lVert#1\right\rVert}
\title{The full relative spectral range of a fractional power:\\an Araki--Horn completion and graph-update envelopes}
\author{Independent research audit}
\date{November 15, 2026}
\begin{document}
\maketitle

\begin{abstract}
Fix the complete generalized spectrum $r$ of an SPD pair $(B,A)$.
For $0<\alpha<1$, the closure of the generalized spectra of
$(B^\alpha,A^\alpha)$ consists exactly of the positive lists whose
logarithms are majorized by $\alpha\log r$. This holds for arbitrary
positive $r$, including signed edits. Necessity is Araki's classical
log-majorization; sufficiency combines the classical Weyl--Horn
triangular realization theorem with a separated-scale Schur-pivot
limit. We give a short self-contained proof of the limit and derive
exact spectral envelopes and strict-count optima. In rank one, we also
characterize exact attainment: precisely the admissible spectra with
pairwise distinct entries above one occur at finite condition number.
The resulting graph
bound uses the small effective-resistance Gram matrix of a batch of
edge additions. All matrix-theory ingredients have substantial prior
art; historical priority for the combined range statement is not claimed.
\end{abstract}

\section{The complete range}
All matrices are real. Spectra and vectors are listed in decreasing
order. For real vectors $x,y\in\mathbb R^n$, write $x\prec y$ when
$\sum_{i=1}^k x_i\le\sum_{i=1}^k y_i$ for $k<n$, with equality
for $k=n$. For positive vectors, $q\prec_{\log}p$ means
$\log q\prec\log p$.

\begin{theorem}[Relative spectral closure]\label{thm:range}
Fix $n\ge1$, $0<\alpha<1$, and a decreasing list
$r_1\ge\cdots\ge r_n>0$. Let $\mathcal R_{\alpha,r}$ consist of
the ordered spectra of
\[
 Q_\alpha(A,B)=A^{-\alpha/2}B^\alpha A^{-\alpha/2},
 \qquad A,B\succ0,
 \qquad \spec(A^{-1/2}BA^{-1/2})=r.
\]
Then
\[
 \overline{\mathcal R_{\alpha,r}}
 =\{q_1\ge\cdots\ge q_n>0:\log q\prec\alpha\log r\}.
\]
The same assertion holds if $A$ is required to be diagonal with entries
$\delta^{n-i}$ and $\delta\downarrow0$ is allowed in the closure.
If $r_i\ge1$ for every $i$, the witnesses satisfy $B\succeq A$ and
$\rank(B-A)=\#\{i:r_i>1\}$, exactly.
\end{theorem}

The last rank statement follows because a congruence of a matrix with
eigenvalues $r_i-1$ has the same inertia. Thus the theorem preserves
the complete relative input spectrum, not merely its determinant.
In the rank-one case $r=(1+\tau,1,\ldots,1)$, majorization is
equivalent to $q_i\ge1$ and $\prod_iq_i=(1+\tau)^\alpha$.

\begin{lemma}[General reverse-Schur fractional limit]\label{lem:pivots}
Let $S\succ0$ be fixed. Define
\[
 a_i=\delta^{n-i},\quad A_\delta=\diag(a_i),\quad
 B_\delta=A_\delta^{1/2}S A_\delta^{1/2}.
\]
Let $d_i>0$ be the reverse Schur pivots of $S$:
\[
 d_i=\frac{\det S[i:n,i:n]}{\det S[i+1:n,i+1:n]},
 \qquad\det S[n+1:n,n+1:n]=1.
\]
Then
\[
 A_\delta^{-\alpha/2}B_\delta^\alpha A_\delta^{-\alpha/2}
 \longrightarrow\diag(d_1^\alpha,\ldots,d_n^\alpha)
 \quad\text{in operator norm}.
\]
Convergence is uniform when $S$ ranges over a compact subset of the
SPD cone.
\end{lemma}

\begin{proof}
Choose constants $0<m\le M$ with $mI\preceq S\preceq MI$.
Write $\nu_1\le\cdots\le\nu_n$ for the eigenvalues of $B_\delta$.
The min--max principle gives
\[
 ma_i\le\nu_i\le Ma_i.
\]
In particular distinct scales imply $\nu_i/\nu_{i+1}\to0$.
For $1\le k\le n$, the elementary symmetric polynomial satisfies
\[
 e_k(B_\delta)
 =\sum_{|I|=k}\det S[I,I]\prod_{i\in I}a_i
 \sim\det S[n-k+1:n,n-k+1:n]\prod_{i=n-k+1}^n a_i.
\]
Every other product has an extra positive power of $\delta$.
Likewise, separation of the $\nu_i$ implies
$e_k(B_\delta)\sim\prod_{i=n-k+1}^n\nu_i$.
Taking successive quotients shows
\[
 \nu_i/a_i\longrightarrow d_i.
\]

Let $v_k$ be orthonormal eigenvectors associated with $\nu_k$,
and write $v_{ik}$ for their coordinates. Positivity of $B_\delta$
and $B_\delta^{-1}$ gives two estimates:
\[
 |v_{ik}|^2\le\frac{(B_\delta)_{ii}}{\nu_k}
 \le C\frac{a_i}{a_k},\qquad
 |v_{ik}|^2\le\nu_k(B_\delta^{-1})_{ii}
 \le C\frac{a_k}{a_i}.
\]
Thus $|v_{ik}|^2\le C\min(a_i/a_k,a_k/a_i)$, and
$|v_{ii}|^2\to1$. Define the nonnegative summands
\[
 t_{ik}=a_i^{-\alpha}\nu_k^\alpha|v_{ik}|^2.
\]
For $k>i$, $t_{ik}=O((a_i/a_k)^{1-\alpha})$; for $k<i$,
$t_{ik}=O((a_k/a_i)^{1+\alpha})$. Therefore
$t_{ii}\to d_i^\alpha$ and $t_{ik}\to0$ whenever $i\ne k$.
These are the diagonal summands of the relative fractional matrix.
For its off-diagonal entries,
\[
 |(Q_\alpha)_{ij}|\le\sum_k\sqrt{t_{ik}t_{jk}}\to0
 \qquad(i\ne j).
\]
For each $k$ at least one factor vanishes and all factors stay bounded.
This proves entrywise, hence norm, convergence in fixed dimension.
On compact SPD sets, $m,M,C$ and the nonzero determinant lower bounds
can all be chosen uniformly, proving the last assertion.
\end{proof}

The successive-Schur asymptotics and eigenvector scale bounds have prior
sources: Usevich--Barthelm\'e~\cite[Theorem 4.2 in the arXiv v2 PDF]{UB} and
Gissler--Auger--Hansen~\cite[Theorem 1]{GAH}. The proof above is included
to check that the power normalization permits precisely $0<\alpha<1$;
the scaling mechanism is not presented as new.

\begin{proof}[Proof of Theorem~\ref{thm:range}]
For $X,Y\succ0$, Araki log-majorization~\cite{Araki} states
\[
 \lambda(X^{\alpha/2}Y^\alpha X^{\alpha/2})
 \prec_{\log}
 \lambda(X^{1/2}YX^{1/2})^\alpha,
 \qquad0<\alpha<1.
\]
Apply this with $X=A^{-1}$ and $Y=B$. This proves necessity.
All admissible $q_i$ lie in the compact interval
$[r_n^\alpha,r_1^\alpha]$, so the necessary set is closed.

Conversely suppose $\log q\prec\alpha\log r$, and put
$d_i=q_i^{1/\alpha}$. The real triangular version of the
Weyl--Horn theorem~\cite{Horn,LiMathias} produces a real upper
triangular matrix $R$ with positive diagonal $\sqrt{d_i}$ and
singular values $\sqrt{r_i}$. Set $S=RR^T$.
It has eigenvalues $r_i$, and its reverse Schur pivots are exactly
$d_i$. Indeed the trailing principal block is the corresponding
trailing triangular factor times its transpose, so its determinant
is the product of those diagonal squares.
Lemma~\ref{lem:pivots} now provides the required witnesses, and
$A_\delta^{-1/2}B_\delta A_\delta^{-1/2}=S$ at every $\delta$.
The rank assertion follows from Sylvester's law of inertia.
\end{proof}

Huang--Tam~\cite[Theorem 3.3]{HT} explicitly characterize squared
Cholesky diagonals on a unitary orbit by multiplicative majorization,
using Kostant convexity~\cite{Kostant}. We use the real Weyl--Horn
construction above to remove any possible real-versus-complex issue.

\section{The exact rank-one attainable spectra}
The closure has a sharp finite-condition-number refinement in rank one.
For $\tau=0$ the relative spectrum is identically one; assume $\tau>0$.

\begin{theorem}[Exact rank-one realization]\label{thm:exact}
Fix $n\ge1$, $0<\alpha<1$, and $\tau>0$.
A decreasing list $q$ is the relative spectrum of
\[
 A^{-\alpha/2}(A+uu^T)^\alpha A^{-\alpha/2},
 \qquad A\succ0,\quad u^TA^{-1}u=\tau,
\]
if and only if for some $1\le m\le n$,
\[
 q_1>\cdots>q_m>1,\qquad
 q_{m+1}=\cdots=q_n=1,\qquad
 \prod_iq_i=(1+\tau)^\alpha.
\]
In particular every repeated value strictly above one occurs only in
the closure. Every admissible list with simple values above one is
attained at finite positive condition number.
\end{theorem}

\begin{proof}
Let $\mathcal K$ be the cyclic subspace of $(A,u)$ and write
$m=\dim\mathcal K$. It and its orthogonal complement reduce both
endpoints. On $\mathcal K^\perp$ the correction $C=Q_\alpha-I$ is zero.
On $\mathcal K$ take an eigenbasis with
$A=\diag(\theta_1<\cdots<\theta_m)$ and $u_i>0$.
The rank-one fractional resolvent identity gives
\[
 C_{ij}=h_i h_j\int_0^\infty
 \frac{w(t)}{(\theta_i+t)(\theta_j+t)}\,dt,
 \quad h_i=u_i\theta_i^{-\alpha/2}>0,
 \quad
 w(t)=\frac{\sin(\pi\alpha)}\pi
 \frac{t^\alpha}{1+\sum_i u_i^2/(\theta_i+t)}>0.
\]
The integral converges and is a Gram matrix. For any row and column
index sets $I,J$ of size $k$, continuous Cauchy--Binet gives
\[
 \frac{\det C[I,J]}{\prod_{i\in I}h_i\prod_{j\in J}h_j}
 =\int_{0<t_1<\cdots<t_k<\infty}
 \det\left[\frac1{\theta_i+t_l}\right]_{i\in I,l\le k}
 \det\left[\frac1{\theta_j+t_l}\right]_{j\in J,l\le k}
 \prod_{l=1}^k w(t_l)\,dt_1\cdots dt_k>0.
\]
The Cauchy determinant formula makes both determinants strictly positive.
One can first truncate to a compact integration interval, then pass
to the limit; the integrand is nonnegative and the entries converge.
Thus the active correction is strictly totally positive, and positive
definite. If its eigenvalues are
$\lambda_1\ge\cdots\ge\lambda_m>0$, every compound matrix has
strictly positive entries. Perron--Frobenius applied to its $k$th
compound implies $\lambda_k>\lambda_{k+1}$ for $k<m$:
otherwise its largest eigenvalue $\lambda_1\cdots\lambda_k$ would be
repeated. This is the classical oscillation-matrix simplicity mechanism.
The determinant lemma gives the required product.

For sufficiency first work in the active dimension $m$.
Put $S_0=\log(1+\tau)$ and $x_i^*=\alpha^{-1}\log q_i$.
Choose a small closed Euclidean ball $K$ centered at $x^*$ in
$\mathcal H=\{x:\sum_i x_i=S_0\}$, of radius $\rho>0$,
contained in $x_1>\cdots>x_m>0$. For $x\in K$ set
\[
 z_i(x)=(e^{x_i}-1)e^{\sum_{j>i}x_j},\quad
 A_\delta=\diag(\delta^{m-i}),\quad
 u_i(x)=\sqrt{(A_\delta)_{ii}z_i(x)}.
\]
The telescoping identity gives $\sum_i z_i(x)=e^{S_0}-1=\tau$.
Write $s_i(x)=\sqrt{z_i(x)}$. The reverse pivots of
$I+s(x)s(x)^T$ are $e^{x_i}$.
Lemma~\ref{lem:pivots}, uniformly for this compact family, gives
\[
 \sup_{x\in K}\norm{Q_\delta(x)-\diag(e^{\alpha x_i})}
 \longrightarrow0.
\]
Consequently the continuous ordered-log-spectrum map
\[
 F_\delta(x)=\alpha^{-1}\log\lambda(Q_\delta(x))
\]
takes values in $\mathcal H$ and converges uniformly to $x$ on $K$.
The affine-space condition follows from the exact determinant identity.
For sufficiently small positive $\delta$,
$T_\delta(x)=x^*+x-F_\delta(x)$ maps $K$ to itself, because
$\norm{T_\delta(x)-x^*}_2<\rho$.
Brouwer gives a fixed point, at which $F_\delta(x)=x^*$.
For $m=1$ use the scalar identity directly. Finally pad inactive
dimensions with any SPD block and zeros in $u$.
\end{proof}

The proof is existential and supplies no useful upper bound on the
condition number of an exact witness. The simplicity statement uses
classical Cauchy-kernel total positivity, while the inverse existence
uses Brouwer; historical priority for this application has not been
established.

\begin{corollary}[Exact rank and simplicity for one graph edge]
\label{cor:graphsimple}
For a positive edge-conductance edit of a connected graph, every
nonzero eigenvalue of the relative fractional correction is positive
and simple. Its rank is
\[
 \#\{\lambda>0:\lambda\in\spec(L),\ P_\lambda b\ne0\},
\]
where $b$ is the edited edge incidence and $P_\lambda$ the spectral
projector of $L$. For an endpoint edge of an $N$-vertex weighted path,
the rank is exactly $N-1$ for every positive edit size.
\end{corollary}
\begin{proof}
Restrict to $\mathbf1^\perp$ and apply the cyclic-subspace argument
above. A weighted path Laplacian is irreducible symmetric tridiagonal,
so its eigenvalues are simple. If an eigenvector $v$ with positive
eigenvalue is orthogonal to an endpoint incidence, then $v_1=v_2$.
The endpoint eigenvector equation gives $\lambda v_1=0$, hence
$v_1=v_2=0$. Recurrence along the path forces $v=0$, a contradiction.
Every positive path eigenspace is therefore seen by that incidence.
\end{proof}

This exact graph obstruction is compatible with low approximate rank.
The path closure theorem permits arbitrarily many eigenvalues tending
to zero. It does not assert exact path realization of each matrix
spectrum with exactly zero padded response eigenvalues.

\section{Exact spectral envelopes and counts}
\begin{corollary}[Every individual spectral envelope is sharp]
\label{cor:envelope}
Under the hypotheses of Theorem~\ref{thm:range}, for $1\le k\le n$,
\[
 \sup\lambda_k(Q_\alpha)
 =\left(\prod_{i=1}^k r_i\right)^{\alpha/k},
 \qquad
 \inf\lambda_{n-k+1}(Q_\alpha)
 =\left(\prod_{i=n-k+1}^n r_i\right)^{\alpha/k}.
\]
The extrema are suprema or infima over actual SPD pairs; attainment
at finite condition number is not asserted.
\end{corollary}
\begin{proof}
If $x=\log\lambda(Q_\alpha)$ and $y=\alpha\log r$, then
$kx_k\le\sum_{i=1}^kx_i\le\sum_{i=1}^ky_i$.
To saturate the first bound in the closure, replace the first $k$
entries of $y$ by their mean and keep all other entries unchanged.
The resulting ordered vector is majorized by $y$.
Average the bottom $k$ entries for the lower bound.
\end{proof}

\begin{corollary}[Exact strict counts for a fixed input spectrum]
\label{cor:count}
For $s>0$,
\[
 \sup\#\{j:\lambda_j(Q_\alpha)>s\}
 =\max\left(\{0\}\cup
 \left\{k\in\{1,\ldots,n\}:
 \alpha\sum_{i=1}^k\log r_i>k\log s\right\}\right).
\]
Likewise
\[
 \sup\#\{j:\lambda_j(Q_\alpha)<s\}
 =\max\left(\{0\}\cup
 \left\{k\in\{1,\ldots,n\}:
 \alpha\sum_{i=n-k+1}^n\log r_i<k\log s\right\}\right).
\]
\end{corollary}
\begin{proof}
The envelope bounds give necessity. If the relevant inequality is
strict, the averaged target spectrum in the preceding proof has
the required $k$ eigenvalues strictly on the desired side. A sufficiently
close finite witness preserves the strict inequalities.
\end{proof}

For a rank-one addition with leverage $\tau$, choose $s=1+\varepsilon$.
For $k\ge1$ the condition is simply
$k\log(1+\varepsilon)<\alpha\log(1+\tau)$.
This recovers the strict ceiling-minus-one formula in the audited
manuscript. For a batch, the first few eigenvalue envelopes can be
strictly stronger than a bound that uses only the determinant.

\section{A graph consequence for a batch of edge additions}
Let $L$ be a connected graph Laplacian on $N$ vertices, and let
\[
 L'=L+UU^T,\qquad
 U=(\sqrt{w_1}b_1,\ldots,\sqrt{w_m}b_m),\quad w_j>0,
\]
where the $b_j$ are edge-incidence vectors.
All matrices in this section act on $\mathbf1^\perp$, of dimension
$d=N-1$. Let $\gamma_1\ge\cdots\ge\gamma_r>0$ be the positive
eigenvalues of the effective-resistance Gram matrix
\[
 G=U^TL^\dagger U,
\]
and pad $\gamma_i=0$ to length $d$. Then the original relative
input eigenvalues are $r_i=1+\gamma_i$.

\begin{corollary}[A complete-Gram spectral envelope]\label{cor:graph}
For
\[
 C=(L^\dagger)^{\alpha/2}(L'^\alpha-L^\alpha)
      (L^\dagger)^{\alpha/2}
\]
on $\mathbf1^\perp$, one has for $1\le k\le d$
\[
 0\le\lambda_k(C)
 \le\exp\left(\frac\alpha k
      \sum_{i=1}^{\min(k,r)}\log(1+\gamma_i)\right)-1.
\]
Consequently a correction $F$ of rank at most $k$ exists with
\[
 0\preceq (L'^\alpha-L^\alpha)-F\preceq\varepsilon L^\alpha
\]
as soon as $k<d$ and
\[
 \alpha\sum_{i=1}^{\min(k+1,r)}\log(1+\gamma_i)
 \le(k+1)\log(1+\varepsilon).
\]
The exact universal matrix worst-case count for this fixed list
$\gamma$ is Corollary~\ref{cor:count} with $s=1+\varepsilon$.
For graph additions the bound always holds; for one edited edge its
unrestricted sharpness on weighted unit-mass paths follows from the
path realization in the audited manuscript.
\end{corollary}
\begin{proof}
The nonunit eigenvalues of $L^{-1/2}L'L^{-1/2}$ on the common support
are $1+\gamma_i$. Apply Corollary~\ref{cor:envelope}.
Truncate $C$ at its largest $k$ eigenvalues, obtaining $C_k$,
and set $F=L^{\alpha/2}C_kL^{\alpha/2}$.
Then $0\preceq C-C_k\preceq\lambda_{k+1}(C)I$, which proves the
energy inequality after congruence.
\end{proof}

This is an existence and obstruction statement for low-rank
preconditioner updates. It does not give a fast method to compute the
spectral truncation. The small Gram matrix supplies a checkable rank
certificate, but computing it requires effective-resistance solves.
No realization of every higher-rank fixed-$\gamma$ extremum by a
prescribed graph class is claimed here.

\section{Priority assessment and relation to the rank-one audit}
The audited rank-one diagonal calculation is correct, including the
strict threshold convention. Its generalization is an exact range
completion of classical inequalities, rather than a stronger necessary
inequality. The necessity is fully contained in Araki; the prescribed
pivot range is fully contained in Weyl--Horn/Kostant; and the underlying
Schur asymptotics are established in the cited perturbation work.
We did not locate this combined range statement verbatim during the
targeted search. That is not evidence that it is historically new.
The graph realization and exact finite rank-one attainable-spectrum
question should be assessed separately from this classical chain.

\begin{thebibliography}{9}
\bibitem{Araki} H. Araki, On an inequality of Lieb and Thirring,
\emph{Lett. Math. Phys.} 19 (1990), 167--170.
\url{https://doi.org/10.1007/BF01045887}.
\bibitem{Horn} A. Horn, On the eigenvalues of a matrix with prescribed
singular values, \emph{Proc. Amer. Math. Soc.} 5 (1954), 4--7.
\url{https://doi.org/10.1090/S0002-9939-1954-0061573-6}.
\bibitem{LiMathias} C.-K. Li and R. Mathias, Construction of matrices
with prescribed singular values and eigenvalues, author manuscript,
October 26, 2004, Sections 3--4.
\url{https://cklixx.people.wm.edu/higham.pdf}.
\bibitem{Kostant} B. Kostant, On convexity, the Weyl group and Iwasawa
decomposition, \emph{Ann. Sci. \`Ecole Norm. Sup.} 6 (1973), 413--455.
\url{https://doi.org/10.24033/asens.1255}.
\bibitem{HT} H. Huang and T.-Y. Tam, Some asymptotic behaviors
associated with matrix decompositions, author manuscript, Theorem 3.3.
\url{https://webhome.auburn.edu/~tamtiny/asymptotic-matrix-lie-6.pdf}.
\bibitem{UB} K. Usevich and S. Barthelm\'e, Computing asymptotic
eigenvectors and eigenvalues of perturbed symmetric matrices,
\emph{SIAM J. Matrix Anal. Appl.} 47 (2026), 429--465.
\url{https://doi.org/10.1137/24M1677460}.
For the cited numbering, see the arXiv v2 PDF:
\url{https://arxiv.org/pdf/2407.17047v2}.
\bibitem{GAH} A. Gissler, A. Auger and N. Hansen, Asymptotic
estimations of a perturbed symmetric eigenproblem (2024).
\url{https://arxiv.org/abs/2403.06983}.
\end{thebibliography}
\end{document}
